\section{}

\paragraph{1076年}\eventanchor{LIAO-1076-REGNAL-YEAR}{LIAO-1076-REGNAL-YEAR}　辽以辽道宗在位第21年作为诏令、赋役与外交文书的纪年框架。账本以《辽史》〈本纪〉、〈营卫志〉、〈百官志〉；《宋史》辽国相关传；宋人使辽记录标明来源路径；该锚点连接编年节点，关于数量、实施和地方后果仍应遵守材料的文体与立场限制。

\paragraph{1077年}\eventanchor{LIAO-1077-REGNAL-YEAR}{LIAO-1077-REGNAL-YEAR}　辽以辽道宗在位第22年作为诏令、赋役与外交文书的纪年框架。账本以《辽史》〈本纪〉、〈营卫志〉、〈百官志〉；《宋史》辽国相关传；宋人使辽记录标明来源路径；该锚点连接编年节点，关于数量、实施和地方后果仍应遵守材料的文体与立场限制。

\paragraph{1078年}\eventanchor{LIAO-1078-REGNAL-YEAR}{LIAO-1078-REGNAL-YEAR}　辽以辽道宗在位第23年作为诏令、赋役与外交文书的纪年框架。账本以《辽史》〈本纪〉、〈营卫志〉、〈百官志〉；《宋史》辽国相关传；宋人使辽记录标明来源路径；该锚点连接编年节点，关于数量、实施和地方后果仍应遵守材料的文体与立场限制。

\paragraph{1079年}\eventanchor{LIAO-1079-REGNAL-YEAR}{LIAO-1079-REGNAL-YEAR}　辽以辽道宗在位第24年作为诏令、赋役与外交文书的纪年框架。账本以《辽史》〈本纪〉、〈营卫志〉、〈百官志〉；《宋史》辽国相关传；宋人使辽记录标明来源路径；该锚点连接编年节点，关于数量、实施和地方后果仍应遵守材料的文体与立场限制。

\paragraph{1080年}\eventanchor{LIAO-1080-REGNAL-YEAR}{LIAO-1080-REGNAL-YEAR}　辽以辽道宗在位第25年作为诏令、赋役与外交文书的纪年框架。账本以《辽史》〈本纪〉、〈营卫志〉、〈百官志〉；《宋史》辽国相关传；宋人使辽记录标明来源路径；该锚点连接编年节点，关于数量、实施和地方后果仍应遵守材料的文体与立场限制。

\paragraph{1081年}\eventanchor{LIAO-1081-REGNAL-YEAR}{LIAO-1081-REGNAL-YEAR}　辽以辽道宗在位第26年作为诏令、赋役与外交文书的纪年框架。账本以《辽史》〈本纪〉、〈营卫志〉、〈百官志〉；《宋史》辽国相关传；宋人使辽记录标明来源路径；该锚点连接编年节点，关于数量、实施和地方后果仍应遵守材料的文体与立场限制。

\paragraph{1082年}\eventanchor{LIAO-1082-REGNAL-YEAR}{LIAO-1082-REGNAL-YEAR}　辽以辽道宗在位第27年作为诏令、赋役与外交文书的纪年框架。账本以《辽史》〈本纪〉、〈营卫志〉、〈百官志〉；《宋史》辽国相关传；宋人使辽记录标明来源路径；该锚点连接编年节点，关于数量、实施和地方后果仍应遵守材料的文体与立场限制。

\paragraph{1083年}\eventanchor{LIAO-1083-REGNAL-YEAR}{LIAO-1083-REGNAL-YEAR}　辽以辽道宗在位第28年作为诏令、赋役与外交文书的纪年框架。账本以《辽史》〈本纪〉、〈营卫志〉、〈百官志〉；《宋史》辽国相关传；宋人使辽记录标明来源路径；该锚点连接编年节点，关于数量、实施和地方后果仍应遵守材料的文体与立场限制。

\paragraph{1084年}\eventanchor{LIAO-1084-REGNAL-YEAR}{LIAO-1084-REGNAL-YEAR}　辽以辽道宗在位第29年作为诏令、赋役与外交文书的纪年框架。账本以《辽史》〈本纪〉、〈营卫志〉、〈百官志〉；《宋史》辽国相关传；宋人使辽记录标明来源路径；该锚点连接编年节点，关于数量、实施和地方后果仍应遵守材料的文体与立场限制。

\paragraph{1085年}\eventanchor{LIAO-1085-REGNAL-YEAR}{LIAO-1085-REGNAL-YEAR}　辽以辽道宗在位第30年作为诏令、赋役与外交文书的纪年框架。账本以《辽史》〈本纪〉、〈营卫志〉、〈百官志〉；《宋史》辽国相关传；宋人使辽记录标明来源路径；该锚点连接编年节点，关于数量、实施和地方后果仍应遵守材料的文体与立场限制。

\paragraph{1086年}\eventanchor{LIAO-1086-REGNAL-YEAR}{LIAO-1086-REGNAL-YEAR}　辽以辽道宗在位第31年作为诏令、赋役与外交文书的纪年框架。账本以《辽史》〈本纪〉、〈营卫志〉、〈百官志〉；《宋史》辽国相关传；宋人使辽记录标明来源路径；该锚点连接编年节点，关于数量、实施和地方后果仍应遵守材料的文体与立场限制。

\paragraph{1087年}\eventanchor{LIAO-1087-REGNAL-YEAR}{LIAO-1087-REGNAL-YEAR}　辽以辽道宗在位第32年作为诏令、赋役与外交文书的纪年框架。账本以《辽史》〈本纪〉、〈营卫志〉、〈百官志〉；《宋史》辽国相关传；宋人使辽记录标明来源路径；该锚点连接编年节点，关于数量、实施和地方后果仍应遵守材料的文体与立场限制。

\paragraph{1088年}\eventanchor{LIAO-1088-REGNAL-YEAR}{LIAO-1088-REGNAL-YEAR}　辽以辽道宗在位第33年作为诏令、赋役与外交文书的纪年框架。账本以《辽史》〈本纪〉、〈营卫志〉、〈百官志〉；《宋史》辽国相关传；宋人使辽记录标明来源路径；该锚点连接编年节点，关于数量、实施和地方后果仍应遵守材料的文体与立场限制。

\paragraph{1089年}\eventanchor{LIAO-1089-REGNAL-YEAR}{LIAO-1089-REGNAL-YEAR}　辽以辽道宗在位第34年作为诏令、赋役与外交文书的纪年框架。账本以《辽史》〈本纪〉、〈营卫志〉、〈百官志〉；《宋史》辽国相关传；宋人使辽记录标明来源路径；该锚点连接编年节点，关于数量、实施和地方后果仍应遵守材料的文体与立场限制。

\paragraph{1090年}\eventanchor{LIAO-1090-REGNAL-YEAR}{LIAO-1090-REGNAL-YEAR}　辽以辽道宗在位第35年作为诏令、赋役与外交文书的纪年框架。账本以《辽史》〈本纪〉、〈营卫志〉、〈百官志〉；《宋史》辽国相关传；宋人使辽记录标明来源路径；该锚点连接编年节点，关于数量、实施和地方后果仍应遵守材料的文体与立场限制。

\paragraph{1091年}\eventanchor{LIAO-1091-REGNAL-YEAR}{LIAO-1091-REGNAL-YEAR}　辽以辽道宗在位第36年作为诏令、赋役与外交文书的纪年框架。账本以《辽史》〈本纪〉、〈营卫志〉、〈百官志〉；《宋史》辽国相关传；宋人使辽记录标明来源路径；该锚点连接编年节点，关于数量、实施和地方后果仍应遵守材料的文体与立场限制。

\paragraph{1092年}\eventanchor{LIAO-1092-REGNAL-YEAR}{LIAO-1092-REGNAL-YEAR}　辽以辽道宗在位第37年作为诏令、赋役与外交文书的纪年框架。账本以《辽史》〈本纪〉、〈营卫志〉、〈百官志〉；《宋史》辽国相关传；宋人使辽记录标明来源路径；该锚点连接编年节点，关于数量、实施和地方后果仍应遵守材料的文体与立场限制。

\paragraph{1093年}\eventanchor{LIAO-1093-REGNAL-YEAR}{LIAO-1093-REGNAL-YEAR}　辽以辽道宗在位第38年作为诏令、赋役与外交文书的纪年框架。账本以《辽史》〈本纪〉、〈营卫志〉、〈百官志〉；《宋史》辽国相关传；宋人使辽记录标明来源路径；该锚点连接编年节点，关于数量、实施和地方后果仍应遵守材料的文体与立场限制。

\paragraph{1094年}\eventanchor{LIAO-1094-EVENT-042}{LIAO-1094-EVENT-042}　辽道宗去世，天祚帝即位。账本以《辽史》〈本纪〉、〈营卫志〉、〈百官志〉；《宋史》辽国相关传；宋人使辽记录标明来源路径；该锚点连接编年节点，关于数量、实施和地方后果仍应遵守材料的文体与立场限制。

\paragraph{1094年}\eventanchor{LIAO-1094-REGNAL-YEAR}{LIAO-1094-REGNAL-YEAR}　辽以辽道宗在位第39年作为诏令、赋役与外交文书的纪年框架。账本以《辽史》〈本纪〉、〈营卫志〉、〈百官志〉；《宋史》辽国相关传；宋人使辽记录标明来源路径；该锚点连接编年节点，关于数量、实施和地方后果仍应遵守材料的文体与立场限制。

\paragraph{1095年}\eventanchor{LIAO-1095-REGNAL-YEAR}{LIAO-1095-REGNAL-YEAR}　辽以辽道宗在位第40年作为诏令、赋役与外交文书的纪年框架。账本以《辽史》〈本纪〉、〈营卫志〉、〈百官志〉；《宋史》辽国相关传；宋人使辽记录标明来源路径；该锚点连接编年节点，关于数量、实施和地方后果仍应遵守材料的文体与立场限制。

\paragraph{1096年}\eventanchor{LIAO-1096-REGNAL-YEAR}{LIAO-1096-REGNAL-YEAR}　辽以辽道宗在位第41年作为诏令、赋役与外交文书的纪年框架。账本以《辽史》〈本纪〉、〈营卫志〉、〈百官志〉；《宋史》辽国相关传；宋人使辽记录标明来源路径；该锚点连接编年节点，关于数量、实施和地方后果仍应遵守材料的文体与立场限制。

\paragraph{1097年}\eventanchor{LIAO-1097-REGNAL-YEAR}{LIAO-1097-REGNAL-YEAR}　辽以辽道宗在位第42年作为诏令、赋役与外交文书的纪年框架。账本以《辽史》〈本纪〉、〈营卫志〉、〈百官志〉；《宋史》辽国相关传；宋人使辽记录标明来源路径；该锚点连接编年节点，关于数量、实施和地方后果仍应遵守材料的文体与立场限制。

\paragraph{1098年}\eventanchor{LIAO-1098-REGNAL-YEAR}{LIAO-1098-REGNAL-YEAR}　辽以辽道宗在位第43年作为诏令、赋役与外交文书的纪年框架。账本以《辽史》〈本纪〉、〈营卫志〉、〈百官志〉；《宋史》辽国相关传；宋人使辽记录标明来源路径；该锚点连接编年节点，关于数量、实施和地方后果仍应遵守材料的文体与立场限制。

\paragraph{1099年}\eventanchor{LIAO-1099-REGNAL-YEAR}{LIAO-1099-REGNAL-YEAR}　辽以辽道宗在位第44年作为诏令、赋役与外交文书的纪年框架。账本以《辽史》〈本纪〉、〈营卫志〉、〈百官志〉；《宋史》辽国相关传；宋人使辽记录标明来源路径；该锚点连接编年节点，关于数量、实施和地方后果仍应遵守材料的文体与立场限制。

\paragraph{1100年}\eventanchor{LIAO-1100-REGNAL-YEAR}{LIAO-1100-REGNAL-YEAR}　辽以辽道宗在位第45年作为诏令、赋役与外交文书的纪年框架。账本以《辽史》〈本纪〉、〈营卫志〉、〈百官志〉；《宋史》辽国相关传；宋人使辽记录标明来源路径；该锚点连接编年节点，关于数量、实施和地方后果仍应遵守材料的文体与立场限制。

\paragraph{1101年}\eventanchor{LIAO-1101-REGNAL-YEAR}{LIAO-1101-REGNAL-YEAR}　辽以辽道宗在位第46年作为诏令、赋役与外交文书的纪年框架。账本以《辽史》〈本纪〉、〈营卫志〉、〈百官志〉；《宋史》辽国相关传；宋人使辽记录标明来源路径；该锚点连接编年节点，关于数量、实施和地方后果仍应遵守材料的文体与立场限制。

\paragraph{1102年}\eventanchor{LIAO-1102-REGNAL-YEAR}{LIAO-1102-REGNAL-YEAR}　辽以辽天祚帝在位第1年作为诏令、赋役与外交文书的纪年框架。账本以《辽史》〈本纪〉、〈营卫志〉、〈百官志〉；《宋史》辽国相关传；宋人使辽记录标明来源路径；该锚点连接编年节点，关于数量、实施和地方后果仍应遵守材料的文体与立场限制。

\paragraph{1103年}\eventanchor{LIAO-1103-REGNAL-YEAR}{LIAO-1103-REGNAL-YEAR}　辽以辽天祚帝在位第2年作为诏令、赋役与外交文书的纪年框架。账本以《辽史》〈本纪〉、〈营卫志〉、〈百官志〉；《宋史》辽国相关传；宋人使辽记录标明来源路径；该锚点连接编年节点，关于数量、实施和地方后果仍应遵守材料的文体与立场限制。

\paragraph{1104年}\eventanchor{LIAO-1104-REGNAL-YEAR}{LIAO-1104-REGNAL-YEAR}　辽以辽天祚帝在位第3年作为诏令、赋役与外交文书的纪年框架。账本以《辽史》〈本纪〉、〈营卫志〉、〈百官志〉；《宋史》辽国相关传；宋人使辽记录标明来源路径；该锚点连接编年节点，关于数量、实施和地方后果仍应遵守材料的文体与立场限制。

\paragraph{1105年}\eventanchor{LIAO-1105-REGNAL-YEAR}{LIAO-1105-REGNAL-YEAR}　辽以辽天祚帝在位第4年作为诏令、赋役与外交文书的纪年框架。账本以《辽史》〈本纪〉、〈营卫志〉、〈百官志〉；《宋史》辽国相关传；宋人使辽记录标明来源路径；该锚点连接编年节点，关于数量、实施和地方后果仍应遵守材料的文体与立场限制。

\paragraph{1106年}\eventanchor{LIAO-1106-REGNAL-YEAR}{LIAO-1106-REGNAL-YEAR}　辽以辽天祚帝在位第5年作为诏令、赋役与外交文书的纪年框架。账本以《辽史》〈本纪〉、〈营卫志〉、〈百官志〉；《宋史》辽国相关传；宋人使辽记录标明来源路径；该锚点连接编年节点，关于数量、实施和地方后果仍应遵守材料的文体与立场限制。

\paragraph{1107年}\eventanchor{LIAO-1107-REGNAL-YEAR}{LIAO-1107-REGNAL-YEAR}　辽以辽天祚帝在位第6年作为诏令、赋役与外交文书的纪年框架。账本以《辽史》〈本纪〉、〈营卫志〉、〈百官志〉；《宋史》辽国相关传；宋人使辽记录标明来源路径；该锚点连接编年节点，关于数量、实施和地方后果仍应遵守材料的文体与立场限制。

\paragraph{1108年}\eventanchor{LIAO-1108-REGNAL-YEAR}{LIAO-1108-REGNAL-YEAR}　辽以辽天祚帝在位第7年作为诏令、赋役与外交文书的纪年框架。账本以《辽史》〈本纪〉、〈营卫志〉、〈百官志〉；《宋史》辽国相关传；宋人使辽记录标明来源路径；该锚点连接编年节点，关于数量、实施和地方后果仍应遵守材料的文体与立场限制。

\paragraph{1109年}\eventanchor{LIAO-1109-REGNAL-YEAR}{LIAO-1109-REGNAL-YEAR}　辽以辽天祚帝在位第8年作为诏令、赋役与外交文书的纪年框架。账本以《辽史》〈本纪〉、〈营卫志〉、〈百官志〉；《宋史》辽国相关传；宋人使辽记录标明来源路径；该锚点连接编年节点，关于数量、实施和地方后果仍应遵守材料的文体与立场限制。

\paragraph{1110年}\eventanchor{LIAO-1110-REGNAL-YEAR}{LIAO-1110-REGNAL-YEAR}　辽以辽天祚帝在位第9年作为诏令、赋役与外交文书的纪年框架。账本以《辽史》〈本纪〉、〈营卫志〉、〈百官志〉；《宋史》辽国相关传；宋人使辽记录标明来源路径；该锚点连接编年节点，关于数量、实施和地方后果仍应遵守材料的文体与立场限制。

\paragraph{1111年}\eventanchor{LIAO-1111-REGNAL-YEAR}{LIAO-1111-REGNAL-YEAR}　辽以辽天祚帝在位第10年作为诏令、赋役与外交文书的纪年框架。账本以《辽史》〈本纪〉、〈营卫志〉、〈百官志〉；《宋史》辽国相关传；宋人使辽记录标明来源路径；该锚点连接编年节点，关于数量、实施和地方后果仍应遵守材料的文体与立场限制。

\paragraph{1112年}\eventanchor{LIAO-1112-REGNAL-YEAR}{LIAO-1112-REGNAL-YEAR}　辽以辽天祚帝在位第11年作为诏令、赋役与外交文书的纪年框架。账本以《辽史》〈本纪〉、〈营卫志〉、〈百官志〉；《宋史》辽国相关传；宋人使辽记录标明来源路径；该锚点连接编年节点，关于数量、实施和地方后果仍应遵守材料的文体与立场限制。

\paragraph{1113年}\eventanchor{LIAO-1113-REGNAL-YEAR}{LIAO-1113-REGNAL-YEAR}　辽以辽天祚帝在位第12年作为诏令、赋役与外交文书的纪年框架。账本以《辽史》〈本纪〉、〈营卫志〉、〈百官志〉；《宋史》辽国相关传；宋人使辽记录标明来源路径；该锚点连接编年节点，关于数量、实施和地方后果仍应遵守材料的文体与立场限制。

\paragraph{1114年}\eventanchor{LIAO-1114-EVENT-043}{LIAO-1114-EVENT-043}　完颜阿骨打起兵反辽，女真军迅速取得初步胜利。账本以《辽史》〈本纪〉、〈营卫志〉、〈百官志〉；《宋史》辽国相关传；宋人使辽记录标明来源路径；该锚点连接编年节点，关于数量、实施和地方后果仍应遵守材料的文体与立场限制。

\paragraph{1114年}\eventanchor{LIAO-1114-REGNAL-YEAR}{LIAO-1114-REGNAL-YEAR}　辽以辽天祚帝在位第13年作为诏令、赋役与外交文书的纪年框架。账本以《辽史》〈本纪〉、〈营卫志〉、〈百官志〉；《宋史》辽国相关传；宋人使辽记录标明来源路径；该锚点连接编年节点，关于数量、实施和地方后果仍应遵守材料的文体与立场限制。

\paragraph{1115年}\eventanchor{LIAO-1115-EVENT-044}{LIAO-1115-EVENT-044}　金建国，辽失去对女真关系的主导权。账本以《辽史》〈本纪〉、〈营卫志〉、〈百官志〉；《宋史》辽国相关传；宋人使辽记录标明来源路径；该锚点连接编年节点，关于数量、实施和地方后果仍应遵守材料的文体与立场限制。

\paragraph{1115年}\eventanchor{LIAO-1115-REGNAL-YEAR}{LIAO-1115-REGNAL-YEAR}　辽以辽天祚帝在位第14年作为诏令、赋役与外交文书的纪年框架。账本以《辽史》〈本纪〉、〈营卫志〉、〈百官志〉；《宋史》辽国相关传；宋人使辽记录标明来源路径；该锚点连接编年节点，关于数量、实施和地方后果仍应遵守材料的文体与立场限制。

\paragraph{1116年}\eventanchor{LIAO-1116-REGNAL-YEAR}{LIAO-1116-REGNAL-YEAR}　辽以辽天祚帝在位第15年作为诏令、赋役与外交文书的纪年框架。账本以《辽史》〈本纪〉、〈营卫志〉、〈百官志〉；《宋史》辽国相关传；宋人使辽记录标明来源路径；该锚点连接编年节点，关于数量、实施和地方后果仍应遵守材料的文体与立场限制。

\paragraph{1117年}\eventanchor{LIAO-1117-REGNAL-YEAR}{LIAO-1117-REGNAL-YEAR}　辽以辽天祚帝在位第16年作为诏令、赋役与外交文书的纪年框架。账本以《辽史》〈本纪〉、〈营卫志〉、〈百官志〉；《宋史》辽国相关传；宋人使辽记录标明来源路径；该锚点连接编年节点，关于数量、实施和地方后果仍应遵守材料的文体与立场限制。

\paragraph{1118年}\eventanchor{LIAO-1118-REGNAL-YEAR}{LIAO-1118-REGNAL-YEAR}　辽以辽天祚帝在位第17年作为诏令、赋役与外交文书的纪年框架。账本以《辽史》〈本纪〉、〈营卫志〉、〈百官志〉；《宋史》辽国相关传；宋人使辽记录标明来源路径；该锚点连接编年节点，关于数量、实施和地方后果仍应遵守材料的文体与立场限制。

\paragraph{1119年}\eventanchor{LIAO-1119-REGNAL-YEAR}{LIAO-1119-REGNAL-YEAR}　辽以辽天祚帝在位第18年作为诏令、赋役与外交文书的纪年框架。账本以《辽史》〈本纪〉、〈营卫志〉、〈百官志〉；《宋史》辽国相关传；宋人使辽记录标明来源路径；该锚点连接编年节点，关于数量、实施和地方后果仍应遵守材料的文体与立场限制。

\paragraph{1120年}\eventanchor{LIAO-1120-EVENT-045}{LIAO-1120-EVENT-045}　宋金结盟共同攻辽，使辽面临南北夹击。账本以《辽史》〈本纪〉、〈营卫志〉、〈百官志〉；《宋史》辽国相关传；宋人使辽记录标明来源路径；该锚点连接编年节点，关于数量、实施和地方后果仍应遵守材料的文体与立场限制。

\paragraph{1120年}\eventanchor{LIAO-1120-REGNAL-YEAR}{LIAO-1120-REGNAL-YEAR}　辽以辽天祚帝在位第19年作为诏令、赋役与外交文书的纪年框架。账本以《辽史》〈本纪〉、〈营卫志〉、〈百官志〉；《宋史》辽国相关传；宋人使辽记录标明来源路径；该锚点连接编年节点，关于数量、实施和地方后果仍应遵守材料的文体与立场限制。

\paragraph{1121年}\eventanchor{LIAO-1121-REGNAL-YEAR}{LIAO-1121-REGNAL-YEAR}　辽以辽天祚帝在位第20年作为诏令、赋役与外交文书的纪年框架。账本以《辽史》〈本纪〉、〈营卫志〉、〈百官志〉；《宋史》辽国相关传；宋人使辽记录标明来源路径；该锚点连接编年节点，关于数量、实施和地方后果仍应遵守材料的文体与立场限制。

\paragraph{1122年}\eventanchor{LIAO-1122-EVENT-046}{LIAO-1122-EVENT-046}　辽南京燕京被金军攻取。账本以《辽史》〈本纪〉、〈营卫志〉、〈百官志〉；《宋史》辽国相关传；宋人使辽记录标明来源路径；该锚点连接编年节点，关于数量、实施和地方后果仍应遵守材料的文体与立场限制。

\paragraph{1122年}\eventanchor{LIAO-1122-REGNAL-YEAR}{LIAO-1122-REGNAL-YEAR}　辽以辽天祚帝在位第21年作为诏令、赋役与外交文书的纪年框架。账本以《辽史》〈本纪〉、〈营卫志〉、〈百官志〉；《宋史》辽国相关传；宋人使辽记录标明来源路径；该锚点连接编年节点，关于数量、实施和地方后果仍应遵守材料的文体与立场限制。

\paragraph{1123年}\eventanchor{LIAO-1123-REGNAL-YEAR}{LIAO-1123-REGNAL-YEAR}　辽以辽天祚帝在位第22年作为诏令、赋役与外交文书的纪年框架。账本以《辽史》〈本纪〉、〈营卫志〉、〈百官志〉；《宋史》辽国相关传；宋人使辽记录标明来源路径；该锚点连接编年节点，关于数量、实施和地方后果仍应遵守材料的文体与立场限制。

\paragraph{1124年}\eventanchor{LIAO-1124-EVENT-047}{LIAO-1124-EVENT-047}　耶律大石率部西行，西辽形成的条件开始出现。账本以《辽史》〈本纪〉、〈营卫志〉、〈百官志〉；《宋史》辽国相关传；宋人使辽记录标明来源路径；该锚点连接编年节点，关于数量、实施和地方后果仍应遵守材料的文体与立场限制。

\paragraph{1124年}\eventanchor{LIAO-1124-REGNAL-YEAR}{LIAO-1124-REGNAL-YEAR}　辽以辽天祚帝在位第23年作为诏令、赋役与外交文书的纪年框架。账本以《辽史》〈本纪〉、〈营卫志〉、〈百官志〉；《宋史》辽国相关传；宋人使辽记录标明来源路径；该锚点连接编年节点，关于数量、实施和地方后果仍应遵守材料的文体与立场限制。

\paragraph{1125年}\eventanchor{LIAO-1125-EVENT-048}{LIAO-1125-EVENT-048}　天祚帝被金军俘获，辽朝统治终结。账本以《辽史》〈本纪〉、〈营卫志〉、〈百官志〉；《宋史》辽国相关传；宋人使辽记录标明来源路径；该锚点连接编年节点，关于数量、实施和地方后果仍应遵守材料的文体与立场限制。

\paragraph{1125年}\eventanchor{LIAO-1125-REGNAL-YEAR}{LIAO-1125-REGNAL-YEAR}　辽以辽天祚帝在位第24年作为诏令、赋役与外交文书的纪年框架。账本以《辽史》〈本纪〉、〈营卫志〉、〈百官志〉；《宋史》辽国相关传；宋人使辽记录标明来源路径；该锚点连接编年节点，关于数量、实施和地方后果仍应遵守材料的文体与立场限制。
